% !TeX TXS-program:compile = txs:///lualatex

\documentclass[a4paper,11pt]{article}
\usepackage[revgoku]{cp-base}
\graphicspath{{./graphics/}}
%variables
\donnees[%
	classe=1\up{ère} 2M2,
	matiere={[SPÉ.MATHS]},
	mois=Février,
	annee=2022,
	typedoc=DM,
	numdoc=07,
	]

%formatage
\author{Pierquet}
\title{\nomfichier}
\hypersetup{pdfauthor={Pierquet},pdftitle={\nomfichier},allbordercolors=white,pdfborder=0 0 0,pdfstartview=FitH}
%divers
\lhead{\entete{\matiere}}
\chead{\entete{\lycee}}
\rhead{\entete{\classe{} - \mois{} \annee}}
%\rhead{\entete{\classe{} - Chapitre }}
\lfoot{\pied{\matiere}}
\cfoot{\logolycee{}}
\rfoot{\pied{\numeropagetot}}
\fancypagestyle{entetedm}{\fancyhead[L]{\entete{\matiere{} À rendre avant le\ldots}}}

\begin{document}

\pagestyle{fancy}

\thispagestyle{entetedm}

\setcounter{numexos}{0}

\part{DM07 - Dérivation locale (facultatif)}

\medskip

\exonum{3}

\medskip

Soit une fonction $f$ définie et dérivable sur $\R$.

Sa courbe représentative {\red $\mathscr{C}_f$} passe par les points {\blue $A(-4\,;\,11)$}, \textcolor{ForestGreen}{$B(2\,;\,4)$} et \textcolor{purple}{$C(6\,;\,2)$}.

Les nombres dérivés de $f$ en $-4$, en $2$ et en $6$ sont respectivement égaux à $-\frac{1}{5}$, à $\frac{1}{5}$ et à $0,465$.

\smallskip

On appelle {\blue $T_A$}, \textcolor{ForestGreen}{$T_B$} et \textcolor{purple}{$T_C$} les tangentes à {\red $\mathscr{C}_f$} respectivement en {\blue $A$}, en \textcolor{ForestGreen}{$B$} et en \textcolor{purple}{$C$}.

\begin{center}
	\tunits{0.45}{0.45}
	\tdefgrille{-10}{14}{2}{2}{-4}{12}{2}{2}
	\begin{tikzpicture}[x=\xunit cm,y=\yunit cm]
		%grilles & axes
		\tgrilles[line width=0.3pt,lightgray!50] ;
		\tgrillep[line width=0.6pt,lightgray!50] ;
		\axestikz* ;
		\axextikz[]{-10,-8,...,12} ;
		\axeytikz[]{-4,-2,...,10} ;
		\clip (\xmin,\ymin) rectangle (\xmax,\ymax) ;
		%splines
		\def\liste{-6/-4/15§-4/11/-0.2§0/4/-0.5§2/4/0.2§6/2/0.465§8/12/8}
		%splines ?
		\splinetikz[liste=\liste,couleur=red,affpoints=false,coeffs=2§3§3§3§2]
		%courbes pgf
%		\draw[line width=1.25pt,red,samples=200,domain=-6:-4] plot(\x,{-0.05*\x*\x*\x+-4.55*\x*\x+-34.2*\x+-56.2}) ;
%		\draw[line width=1.25pt,red,samples=200,domain=-4:0] plot(\x,{0.175*\x*\x*\x+1.0125*\x*\x+-0.5*\x+4.0}) ;
%		\draw[line width=1.25pt,red,samples=200,domain=0:2] plot(\x,{-0.075*\x*\x*\x+0.4*\x*\x+-0.5*\x+4.0}) ;
%		\draw[line width=1.25pt,red,samples=200,domain=2:6] plot(\x,{0.1041*\x*\x*\x+-1.2156*\x*\x+3.8138*\x+0.4025}) ;
%		\draw[line width=1.25pt,red,samples=200,domain=6:8] plot(\x,{-0.6338*\x*\x*\x+15.1925*\x*\x+-113.4*\x+272.36}) ;
		%tangentes
		\draw[line width=1.25pt,densely dashed,blue,domain=\xmin:\xmax,samples=2] plot (\x,{-0.2*(\x+4)+11}) ;
		\draw[line width=1.25pt,densely dashed,ForestGreen,domain=\xmin:\xmax,samples=2] plot (\x,{0.2*(\x-2)+4}) ;
		\draw[line width=1.25pt,densely dashed,purple,domain=\xmin:\xmax,samples=2] plot (\x,{0.465*(\x-6)+2}) ;
		%points
		\foreach \Point/\Nom/\Couleur/\Position in {(-4,11)/A/blue/above right,(2,4)/B/ForestGreen/above left,(6,2)/C/purple/above left} \draw[\Couleur,fill=\Couleur] \Point circle[radius=2pt] node[\Couleur,\Position] {\ppoint{\Nom}} ;
		%labels
		\draw (-7,11) node[blue] {\large $\mathsf{T_A}$} ;
		\draw (-9,1) node[ForestGreen] {\large $\mathsf{T_B}$} ;
		\draw (9,2.5) node[purple] {\large $\mathsf{T_C}$} ;
		\draw (8.5,10.75) node[red] {\Large $\mathscr{C}_f$} ;
	\end{tikzpicture}
\end{center}

\begin{enumerate}
	\item Déterminer l'équation réduite de chacune des tangentes {\blue $T_A$}, \textcolor{ForestGreen}{$T_B$} et \textcolor{purple}{$T_C$}.
	\item 
	\begin{enumerate}
		\item Déterminer les coordonnées du point d'intersection des droites {\blue $T_A$} et \textcolor{ForestGreen}{$T_B$}.
		\item Les trois tangentes {\blue $T_A$}, \textcolor{ForestGreen}{$T_B$} et \textcolor{purple}{$T_C$} sont-elles concourantes ? Justifier.
	\end{enumerate}
\end{enumerate}

\medskip

\exonum{3}

\medskip

On donne une capture d'écran obtenue à l'aide d'un logiciel de calcul formel :

\begin{center}
	\begin{tikzpicture}[x=1cm,y=1cm,line width=1pt]
		\paramCF[larg=14,esplg=2pt,titre=true]
		\ligneCF{\textsf{f(x):=sqrt(x+8)}}{$\mathsf{x \mapsto \sqrt{x+8}}$}
		\ligneCF{\textsf{limite((f(1+h)-f(1))/h,h,0)}}{$\mathsf{\tfrac16}$}
		\ligneCF{\textsf{tangente(f,1)}}{$\mathsf{y=\tfrac16x+\tfrac{17}{6}}$}
	\end{tikzpicture}
\end{center}

\begin{enumerate}
	\item Que permet de déterminer la commande saisie en \cxcas{L2} ?
	\item 
	\begin{enumerate}
		\item Que permet de déterminer la commande saisie en \cxcas{L3} ?
		\item Justifier, par calculs, le résultat obtenu à la ligne \cxcas{L3}.
	\end{enumerate}
	\item[Bonus :] Justifier, par calculs, le résultat obtenu à la ligne \cxcas{L2}.
\end{enumerate}

\end{document}